% Part: first-order-logic
% Chapter: proof-systems
% Section: natural-deduction

\documentclass[../../../include/open-logic-section]{subfiles}

\begin{document}

\iftag{FOL}
      {\olfileid{fol}{prf}{ntd}}
      {\olfileid{pl}{prf}{ntd}}

\olsection{સ્વાભાવિક નિગમન}

સ્વાભાવિક નિગમન વાસ્તવિક તર્કપ્રક્રિયાનું (ખાસ કરીને ગણિતશાસ્ત્રીઓ
વાપરે છે તેવી સુવ્યવસ્થિત તર્કપ્રક્રિયાનું) પ્રતિબિંબ પાડવા રચાયેલું
!!a{derivation} તંત્ર છે. વાસ્તવિક તર્કપ્રક્રિયા ઘણી ``સ્વાભાવિક''
ભાતો મુજબ આગળ વધે છે. દાખલા તરીકે, કિસ્સાવાર સાબિતીમાં કોઈ વિયોજક
આધારવિધાન પરથી ફલિતવિધાન સ્થાપવા માટે તે બંને વિયોજિત ઘટકોમાંથી
ફલિત થાય છે તેમ અલગ અલગ સ્થાપીએ છીએ. પરોક્ષ સાબિતીમાં કોઈ
ફલિતવિધાનનો નિષેધ વ્યાઘાત તરફ દોરી જાય છે તેમ બતાવીને ફલિતવિધાન
સ્થાપીએ છીએ. શરતી સાબિતીમાં પૂર્વવર્તીમાંથી અનુવર્તી ફલિત થાય છે તેમ
બતાવીને ``જો \dots તો \dots'' સ્વરૂપનો પ્રેરણરૂપ દાવો સ્થાપીએ છીએ.
સ્વાભાવિક નિગમન આવાં કેટલાંક સ્વાભાવિક અનુમાનોનું ઔપચારિક રૂપ છે.
દરેક તાર્કિક સંયોજક અને પરિમાણકને બે નિયમો—એક પરિચય-નિયમ અને એક
નિવારણ-નિયમ—હોય છે, અને દરેક નિયમ આવી કોઈ એક સ્વાભાવિક અનુમાન-ભાતને
અનુરૂપ હોય છે. દાખલા તરીકે, $\Intro{\lif}$ શરતી સાબિતીને અને
$\Elim{\lor}$ કિસ્સાવાર સાબિતીને અનુરૂપ છે. ખાસ સરળ નિયમ
$\Elim{\land}$ છે, જે $!A \land !B$ પરથી~$!A$ (અથવા $!B$)નું અનુમાન
કરવાની છૂટ આપે છે.

સ્વાભાવિક નિગમનને બીજાં !!{derivation} તંત્રોથી જુદું પાડતું એક
લક્ષણ તેનો ધારણાઓનો ઉપયોગ છે. સ્વાભાવિક નિગમનમાં !!^a{derivation} એ
!!{formula}sનું વૃક્ષ છે. વૃક્ષના મૂળે એક !!{formula} હોય છે અને
!!{formula}sના વૃક્ષનાં ``પર્ણો'' એવાં !!{formula}s છે જેમાંથી
ફલિતવિધાન નિષ્પન્ન થાય છે. સ્વાભાવિક નિગમનમાં કેટલાંક
પર્ણ-!!{formula}s !!{derivation}ની અંદર
ભૂમિકા ભજવે છે, પરંતુ !!{derivation} ફલિતવિધાન સુધી પહોંચે ત્યાં સુધીમાં
તેમનો ઉપયોગ પૂરો થઈ જાય છે. વાસ્તવિક તર્કપ્રક્રિયામાં થોડા સમય માટે જ
અસરમાં રહેતી પરિકલ્પનાઓ દાખલ કરવાની પ્રથાને આ અનુરૂપ છે. દાખલા તરીકે,
કિસ્સાવાર સાબિતીમાં આપણે દરેક વિયોજિત ઘટક સત્ય છે એમ ધારીએ છીએ; શરતી
સાબિતીમાં પૂર્વવર્તી સત્ય છે એમ ધારીએ છીએ; અને પરોક્ષ સાબિતીમાં
ફલિતવિધાનનો નિષેધ સત્ય છે એમ ધારીએ છીએ. વચલું પગલું સ્થાપવા માટે આવી
પરિકલ્પનાત્મક ધારણાઓ દાખલ કરી પછી તેમને નિવૃત્ત કરવાની રીત સ્વાભાવિક
નિગમનની ખાસ નિશાની છે. સ્વાભાવિક નિગમનની !!{derivation}નાં પર્ણો પરનાં
સૂત્રોને ધારણાઓ કહે છે, અને કેટલાંક અનુમાન-નિયમો તેમને
``!!{discharge}'' કરી શકે છે. દાખલા તરીકે, કેટલીક એવી ધારણાઓમાંથી
આપણી પાસે~$!B$ની !!a{derivation} હોય જેમાં~$!A$નો સમાવેશ થતો હોય, તો
નિયમ $\Intro{\lif}$ વડે~$!A \lif !B$નું અનુમાન કરી શકાય છે અને
સ્વરૂપ~$!A$ની કોઈપણ ધારણાને નિવૃત્ત કરી શકાય છે. (કઈ ધારણાઓને
કયા અનુમાન વખતે નિવૃત્ત કરવામાં આવે છે તેની નોંધ રાખવા, અનુમાનને અને
તે નિવૃત્ત કરે છે તે ધારણાઓને એક સંખ્યા આપીએ છીએ.) !!{derivation}ના
અંતે !!{undischarged} રહેતી ધારણાઓ ફલિતવિધાનના સત્ય માટે એકસાથે પૂરતી
હોય છે; તેથી !!a{derivation} સ્થાપે છે કે તેની !!{undischarged}
ધારણાઓમાંથી તેનું ફલિતવિધાન અર્થાનુસારી રીતે ફલિત થાય છે.

સ્વાભાવિક નિગમન પર આધારિત સંબંધ $\Gamma \Proves !A$ ત્યારે અને માત્ર
ત્યારે જ સધાય છે જ્યારે એવી !!a{derivation} હોય જેમાં વૃક્ષનું છેલ્લું
!!{sentence} $!A$ હોય અને દરેક !!{undischarged} પર્ણ~$\Gamma$માં હોય.
સ્વાભાવિક નિગમનમાં $!A$ ત્યારે અને માત્ર ત્યારે જ પ્રમેય છે જ્યારે એવી
!!a{derivation} હોય જેમાં $!A$ છેલ્લું !!{sentence} હોય અને બધી ધારણાઓ
!!{discharged} હોય. દાખલા તરીકે, નીચેની !!a{derivation} બતાવે છે કે
$\Proves (!A \land !B) \lif !A$:
\begin{prooftree}
  \AxiomC{$\Discharge{!A \land !B}{1}$}
  \RightLabel{\Elim{\land}}
  \UnaryInfC{$!A$}
  \DischargeRule{\Intro{\lif}}{1}
  \UnaryInfC{$(!A \land !B) \lif !A$}
\end{prooftree}
નિશાની~$1$ બતાવે છે કે ધારણા $!A \land !B$ને \Intro{\lif}
અનુમાન વખતે !!{discharged} કરવામાં આવે છે.

સ્વાભાવિક નિગમનમાં $\Gamma \Proves \lfalse$ હોય તો અને માત્ર તો ગણ
$\Gamma$ અસુસંગત છે. નિયમ \FalseInt{} વડે અસુસંગત ગણમાંથી કોઈપણ
!!{sentence} !!{derive}d કરી શકાય છે.

સ્વાભાવિક નિગમનનાં તંત્રો ગેરહાર્ડ ગેન્ટ્ઝન અને સ્તાનિસ્વાવ
યાશ્કોવ્સ્કીએ 1930ના દાયકામાં વિકસાવ્યાં હતાં; પછી ડાગ પ્રાવિટ્ઝ અને
ફ્રેડરિક ફિચે તેમનો વધુ વિકાસ કર્યો. તેનાં અનુમાનો સાબિતીની સ્વાભાવિક
રીતોનું પ્રતિબિંબ પાડતાં હોવાથી તત્ત્વચિંતકો તેને પસંદ કરે છે. ફિચે
વિકસાવેલાં રૂપો તર્કશાસ્ત્રનાં પ્રારંભિક પાઠ્યપુસ્તકોમાં ઘણી વાર
વપરાય છે. તર્કશાસ્ત્રના તત્ત્વચિંતનમાં ક્યારેક એવું માનવામાં આવ્યું છે
કે સ્વાભાવિક નિગમનના નિયમો તાર્કિક કારકોના અર્થ આપે છે
(``સાબિતી-સૈદ્ધાંતિક અર્થવિચાર'').

\end{document}
